\section{创业复盘：低谷、爆发和下一阶段}

本章把访谈的情绪线转成创业复盘。陈冕的“好爽”来自阶段性胜利，但下一阶段要回答的问题更硬：团队如何组织，现金流如何稳住，产品如何留存，商业模式如何成立，设计师专业心智如何占住。AI 应用创业经常会在爆发后进入第二次考试：用户来了以后，系统能不能承载？收入来了以后，组织能不能扩张？竞品来了以后，壁垒还在不在？

\begin{knowledgebox}{爆发后的第二次考试}
\begin{tabularx}{\textwidth}{p{0.20\textwidth}p{0.34\textwidth}X}
\toprule
考试 & 关键问题 & 通过信号 \\
\midrule
产品考试 & 用户是否反复使用同一工作流？ & 留存、项目数、协作深度。 \\
商业考试 & 是否有稳定付费理由？ & 订阅、团队合同、付费转化。 \\
组织考试 & 团队是否能从战斗状态进入系统化？ & 研发节奏、客服、交付、财务纪律。 \\
竞争考试 & 大模型和竞品追上后还剩什么？ & 数据、品牌、资产和生态壁垒。 \\
\bottomrule
\end{tabularx}
\end{knowledgebox}

\begin{warningbox}{爆发之后不要只追更大声量}
AI 应用出圈很容易让团队继续追传播，但真正危险的是基础能力跟不上：稳定性、成本、版权、协作和售后如果没有跟上，声量会反过来放大产品问题。
\end{warningbox}

\subsection{创业者情绪如何转成组织能力}

战斗型 CEO 能穿过低谷，但公司不能长期只靠创始人的情绪燃烧。下一阶段需要把创始人的判断转成组织机制：产品路线图、用户反馈系统、财务纪律、模型评测、版权流程和客户成功。否则，一次爆发可能只是短暂高峰。

\begin{knowledgebox}{从个人能量到组织机制}
\begin{tabularx}{\textwidth}{p{0.22\textwidth}p{0.34\textwidth}X}
\toprule
个人能量 & 组织化之后 & 作用 \\
\midrule
产品直觉 & 用户研究和评测体系 & 避免只靠创始人拍板。 \\
战斗状态 & 节奏管理和优先级机制 & 防止团队长期透支。 \\
融资压力 & 财务模型和 runway 管理 & 避免再次现金见底。 \\
用户兴奋 & 客户成功和社区运营 & 把情绪转成留存。 \\
\bottomrule
\end{tabularx}
\end{knowledgebox}

\subsection{本章小结}

Lovart 的下一阶段不是证明“能火”，而是证明“能稳定成为设计工作台”。这要求从创始人情绪转向组织机制，从 Magic Moment 转向可重复交付。

